\section{A15.1-P: paired-source identities P1 and P2}
\label{sec:p12-P1P2}

We record two unconditional algebraic identities of the localised-hub
kernel operator $L$ in the three-shift chamber, valid on every
defect strip.  They follow from six horizon-legal raw source
equations and eliminate the tail $H$ against the reflected value $l$
and against the lower-half $D$-form.

Notation.  For a candidate $h \in L^{2}(R, S)$ with $L h = 0$,
write $H(t) := h(T + t)$ for the tail and $l(x) := h(T - x)$ for the
reflection through the anchor $T$.  Set $\Delta := p^{2} - q^{2}$.

\begin{lemma}[Six horizon-legal raw source equations]
\label{lem:p12-P-six}
Let $2a < T_{0} < c$, $\varepsilon := T_{0} - T \in (0, \varepsilon_{\max})$
with $\varepsilon_{\max} = \tfrac12\log(5/4)$, $0 < x < \varepsilon$.
The six source positions
$a - x,\ a + x,\ b - x,\ b + x,\ T - x,\ T + x$
lie in $(0, T_{0})$, and the raw operator equations at these sources
read, using the anti-reflected extension $h(-y) = -h(y)$ and
zero-extension outside $(R, S)$:
\begin{align*}
A^{-}(u = a - x): & \quad -p h(x) - p l(x) - r h(d + x) - r H(d - x) - q h(a + x) = 0,\\
A^{+}(u = a + x): & \quad p h(x) - p H(x) - r h(d - x) - q h(a - x) = 0,\\
B^{-}(u = b - x): & \quad p h(d - x) - p H(d - x) - r h(x) - q h(e + x) = 0,\\
B^{+}(u = b + x): & \quad p h(d + x) + r h(x) - q h(e - x) = 0,\\
T^{-}(u = T - x): & \quad p h(a - x) + r h(e - x) - q h(x) = 0,\\
T^{+}(u = T + x): & \quad p h(a + x) + r h(e + x) + q h(x) = 0.
\end{align*}
\end{lemma}

\begin{proof}
Horizon-legality per source (each strictly less than $T_{0}$; the
chain "$<$ some constant $<c$" is \emph{not} used, since $T_{0} < c$):
\begin{itemize}
\item $a - x < a < T < T_{0}$;
\item $a + x < a + \varepsilon_{\max} < T < T_{0}$
  (since $a + \varepsilon_{\max} = a + \tfrac12\log(5/4)
  = \tfrac12\log(5) - \tfrac12\log 2 = c - a < 2a = T$,
  equivalent to $2a > c/2$, i.e.\ $2\log 2 > \tfrac12\log 5$,
  elementary since $16 > 5$);
\item $b - x < b < T$ (since $2a > b$, i.e.\ $4 > 3$);
\item $b + x < b + \varepsilon_{\max} < T$ (since
  $b + \varepsilon_{\max} = \tfrac12(\log 3 + \log(5/4))
  = \tfrac12\log(15/4)$, and $2a = \log 2$;
  $\log 2 > \tfrac12\log(15/4)$ iff $4 > 15/4$, i.e.\ $16 > 15$,
  elementary);
\item $T - x < T < T_{0}$;
\item $T + x < T + \varepsilon = T_{0}$.
\end{itemize}
Operator convention: $Lh(u) = \sum_{s} k_{s}[h(u - s) - h(u + s)]$
with $(k_{a}, k_{b}, k_{2a}) = (p, r, q)$.  Applying at each source
and dropping terms whose $h$-argument lies below $0$ (anti-reflected
to their positive image) or above $S = T + \sigma$ (outside support,
zero) yields the six equations exactly.  Anti-reflection at zero
uses $h(-y) = -h(y)$, which is the natural convention for an
antisymmetric shift operator on the symmetric extension.
\end{proof}

\begin{theorem}[Identity P1]\label{thm:p12-P1}
Under the hypotheses of Lemma~\ref{lem:p12-P-six}, for every
$0 < x < \varepsilon$:
\[
\boxed{\; H(x) + l(x) + \frac{2 r}{p}\, H(d - x) = 0.\;}
\]
\end{theorem}

\begin{proof}
Apply the row-multiplier vector
$\bigl(-\tfrac{1}{p}, -\tfrac{1}{p}, -\tfrac{r}{p^{2}},
-\tfrac{r}{p^{2}}, -\tfrac{q}{p^{2}}, -\tfrac{q}{p^{2}}\bigr)$
to $(A^{-}, A^{+}, B^{-}, B^{+}, T^{-}, T^{+})$; all
$h$-terms cancel identically, leaving $H(x) + l(x) + 2r/p \cdot H(d - x)$.
\end{proof}

\begin{theorem}[Identity P2]\label{thm:p12-P2}
Under the same hypotheses, for every $0 < x < \varepsilon$:
\[
\boxed{\; H(x) - l(x) - \frac{2 D(x)}{p^{2}} = 0,\qquad
D(x) := (p^{2} - q^{2} - r^{2})\, h(x) + q r\,[h(e - x) - h(e + x)].\;}
\]
\end{theorem}

\begin{proof}
Row-multiplier vector
$\bigl(+\tfrac{1}{p}, -\tfrac{1}{p}, -\tfrac{r}{p^{2}},
+\tfrac{r}{p^{2}}, -\tfrac{q}{p^{2}}, +\tfrac{q}{p^{2}}\bigr)$
against $(A^{-}, A^{+}, B^{-}, B^{+}, T^{-}, T^{+})$.  The
$h$-argument book-keeping cancels every visible $h$-value except
the $D$-form.
\end{proof}

\begin{corollary}[Case-B non-degeneracy]
\label{cor:p12-P1-caseB}
For $\sigma > d/2$ and $x \in I := (d - \sigma, \sigma)$, both
$H(x)$ and $H(d - x)$ are live tail values.  Applying P1 at $x$ and
at $d - x$ gives the system
\[
\begin{pmatrix} 1 & c_{0} \\ c_{0} & 1 \end{pmatrix}
\begin{pmatrix} H(x) \\ H(d - x) \end{pmatrix}
= -\begin{pmatrix} l(x) \\ l(d - x) \end{pmatrix},
\qquad c_{0} := 2r / p.
\]
$c_{0} > 1$ elementarily ($4 r^{2} > p^{2} \iff \log 3 / \log 2 >
3\sqrt{3}/(8\sqrt{2}) < 1$, which follows from $\log 3 > \log 2$);
determinant $1 - c_{0}^{2} \ne 0$.  Hence $H(x), H(d - x)$ are
uniquely determined by $l(x), l(d - x)$.
\end{corollary}

\begin{corollary}[Sub-case $\sigma \le d/2$]
\label{cor:p12-P1-subhalf}
Under $\sigma \le d/2$, for every $x \in (0, \sigma)$:
$d - x \ge d/2 \ge \sigma$, so $H(d - x) = 0$.  P1 gives
$H(x) = -l(x)$.  For
$x \in (\sigma, \min(\varepsilon, d - \sigma))$: $H(x) = 0$ and
$H(d - x) = 0$, so P1 forces $l(x) = 0$.  Consequently $h$ vanishes
on the horizon-generated interval
$(T - \min(\varepsilon, d - \sigma), T - \sigma) \subset (a, T)$.
\end{corollary}
